% GENERATED FILE. Edit canonical lessons, then run npm run generate:books.
% canonical-source-hash: fnv1a64:13572fe97679b8fc
% canonical-lessons: JA-C473-nikomu, JA-C473-tokasu, JA-C473-kogasu, JA-C473-fukuramu, JA-C473-samasu, JA-R473-review-468-473

\chapter{Simmering and Melting}
\label{ch:ja-simmering-and-melting}
\hlchaptermodality{473}

\begin{chapteropening}
\textbf{By the end of this chapter:} \emph{I can say how something is cooked.}

Complete the last lesson of chapter 473: I can say how something is cooked.
\end{chapteropening}

\section[\texorpdfstring{\ja{にこむ} (nikomu)}{にこむ (nikomu)}]{\ja{にこむ} (nikomu) — to stew; to simmer}

\label{lesson:JA-C473-nikomu}



\begin{quote}
Before the new one: say the Japanese for wheat flour, then the Japanese for soybeans.
\end{quote}

\subsection*{You'll want to know: \ja{にこむ}}
\textbf{\ja{にこむ}} — \emph{nikomu} — \textquotedblleft{}to stew; to simmer\textquotedblright{}.

\begin{quote}
\textbf{\ja{やさいを} \ja{ながく} \ja{にこんで} \ja{ください。}} — \emph{yasai o nagaku nikonde kudasai} — \textquotedblleft{}Please simmer the vegetables for a long time.\textquotedblright{}
\end{quote}

Polite \textbf{\ja{にこみます}}, \emph{nikomimasu}; a godan verb. Past (polite): \textbf{\ja{ははが} \ja{にくを} \ja{にこみました。}}, \emph{haha ga niku o nikomimashita}.

\begin{cousinweb}
\textbf{\ja{にこむ}} (\emph{nikomu}) is a native Japanese compound (wago), inherited rather than borrowed from Chinese: \emph{ni}, \textquotedblleft{}boiling (niru)\textquotedblright{} and \emph{komu}, \textquotedblleft{}going in\textquotedblright{} (JMdict, KANJIDIC2).
\end{cousinweb}

\begin{grammarlens}[title={What you've built}]
One more word for describing what you lived and felt.
\end{grammarlens}

\subsection*{Your turn}
\begin{itemize}
\raggedright
  \item \emph{Say it:} \emph{nikomu}
  \item \emph{Say it:} \emph{nikomu}, once more
  \item \emph{Say it:} \emph{yasai o nagaku nikonde kudasai} — Please simmer the vegetables for a long time.
  \item \emph{Recall:} say the Japanese for wheat flour, then the Japanese for soybeans, then say \emph{nikomu} again
\end{itemize}

\begin{tcolorbox}[breakable,colback=teal!4,colframe=teal!35!black,title={Before you move on}]
What does \ja{にこむ} mean? (\textquotedblleft{}To stew; to simmer\textquotedblright{}.) Say it once more.
\end{tcolorbox}
\section[\texorpdfstring{\ja{とかす} (tokasu)}{とかす (tokasu)}]{\ja{とかす} (tokasu) — to melt something; to dissolve}

\label{lesson:JA-C473-tokasu}



\begin{quote}
Before the new one: say the Japanese for soybeans, then the Japanese for to stew.
\end{quote}

\subsection*{You'll want to know: \ja{とかす}}
\textbf{\ja{とかす}} — \emph{tokasu} — \textquotedblleft{}to melt something; to dissolve\textquotedblright{}.

\begin{quote}
\textbf{\ja{なべで} \ja{さとうを} \ja{とかしました。}} — \emph{nabe de satō o tokashimashita} — \textquotedblleft{}I melted the sugar in a pot.\textquotedblright{}
\end{quote}

Polite \textbf{\ja{とかします}}, \emph{tokashimasu}; a godan verb. Past (polite): \textbf{\ja{ちちが} \ja{こおりを} \ja{とかしました。}}, \emph{chichi ga kōri o tokashimashita}.

\begin{cousinweb}
\textbf{\ja{とかす}} (\emph{tokasu}) is a native Japanese word (wago), inherited rather than borrowed from Chinese: the transitive partner of \emph{tokeru}, \textquotedblleft{}to melt\textquotedblright{}, and its kanji was attached later as a kun'yomi (JMdict, KANJIDIC2).
\end{cousinweb}

\begin{grammarlens}[title={What you've built}]
One more word for describing what you lived and felt.
\end{grammarlens}

\subsection*{Your turn}
\begin{itemize}
\raggedright
  \item \emph{Say it:} \emph{tokasu}
  \item \emph{Say it:} \emph{tokasu}, once more
  \item \emph{Say it:} \emph{nabe de satō o tokashimashita} — I melted the sugar in a pot.
  \item \emph{Recall:} say the Japanese for soybeans, then the Japanese for to stew, then say \emph{tokasu} again
\end{itemize}

\begin{tcolorbox}[breakable,colback=teal!4,colframe=teal!35!black,title={Before you move on}]
What does \ja{とかす} mean? (\textquotedblleft{}To melt something; to dissolve\textquotedblright{}.) Say it once more.
\end{tcolorbox}
\section[\texorpdfstring{\ja{こがす} (kogasu)}{こがす (kogasu)}]{\ja{こがす} (kogasu) — to scorch; to burn (food)}

\label{lesson:JA-C473-kogasu}



\begin{quote}
Before the new one: say the Japanese for to stew, then the Japanese for to melt something.
\end{quote}

\subsection*{You'll want to know: \ja{こがす}}
\textbf{\ja{こがす}} — \emph{kogasu} — \textquotedblleft{}to scorch; to burn (food)\textquotedblright{}.

\begin{quote}
\textbf{\ja{なべの} \ja{ごはんを} \ja{こがしました。}} — \emph{nabe no gohan o kogashimashita} — \textquotedblleft{}I scorched the rice in the pot.\textquotedblright{}
\end{quote}

Polite \textbf{\ja{こがします}}, \emph{kogashimasu}; a godan verb. Past (polite): \textbf{\ja{あにが} \ja{さかなを} \ja{こがしました。}}, \emph{ani ga sakana o kogashimashita}.

\begin{cousinweb}
\textbf{\ja{こがす}} (\emph{kogasu}) is a native Japanese word (wago), inherited rather than borrowed from Chinese: the transitive partner of \emph{kogeru}, \textquotedblleft{}to get burnt\textquotedblright{}, and its kanji was attached later as a kun'yomi (JMdict, KANJIDIC2).
\end{cousinweb}

\begin{grammarlens}[title={What you've built}]
One more word for describing what you lived and felt.
\end{grammarlens}

\subsection*{Your turn}
\begin{itemize}
\raggedright
  \item \emph{Say it:} \emph{kogasu}
  \item \emph{Say it:} \emph{kogasu}, once more
  \item \emph{Say it:} \emph{nabe no gohan o kogashimashita} — I scorched the rice in the pot.
  \item \emph{Recall:} say the Japanese for to stew, then the Japanese for to melt something, then say \emph{kogasu} again
\end{itemize}

\begin{tcolorbox}[breakable,colback=teal!4,colframe=teal!35!black,title={Before you move on}]
What does \ja{こがす} mean? (\textquotedblleft{}To scorch; to burn (food)\textquotedblright{}.) Say it once more.
\end{tcolorbox}
\section[\texorpdfstring{\ja{ふくらむ} (fukuramu)}{ふくらむ (fukuramu)}]{\ja{ふくらむ} (fukuramu) — to swell; to puff up}

\label{lesson:JA-C473-fukuramu}



\begin{quote}
Before the new one: say the Japanese for to melt something, then the Japanese for to scorch.
\end{quote}

\subsection*{You'll want to know: \ja{ふくらむ}}
\textbf{\ja{ふくらむ}} — \emph{fukuramu} — \textquotedblleft{}to swell; to puff up\textquotedblright{}.

\begin{quote}
\textbf{\ja{もちが} \ja{おおきく} \ja{ふくらみました。}} — \emph{mochi ga ōkiku fukuramimashita} — \textquotedblleft{}The rice cake puffed up big.\textquotedblright{}
\end{quote}

Polite \textbf{\ja{ふくらみます}}, \emph{fukuramimasu}; a godan verb. Past (polite): \textbf{\ja{もちが} \ja{ふくらみました。}}, \emph{mochi ga fukuramimashita}.

\begin{cousinweb}
\textbf{\ja{ふくらむ}} (\emph{fukuramu}) is a native Japanese word (wago), inherited rather than borrowed from Chinese; its kanji spelling was attached to it later (JMdict).
\end{cousinweb}

\begin{grammarlens}[title={What you've built}]
One more word for describing what you lived and felt.
\end{grammarlens}

\subsection*{Your turn}
\begin{itemize}
\raggedright
  \item \emph{Say it:} \emph{fukuramu}
  \item \emph{Say it:} \emph{fukuramu}, once more
  \item \emph{Say it:} \emph{mochi ga ōkiku fukuramimashita} — The rice cake puffed up big.
  \item \emph{Recall:} say the Japanese for to melt something, then the Japanese for to scorch, then say \emph{fukuramu} again
\end{itemize}

\begin{tcolorbox}[breakable,colback=teal!4,colframe=teal!35!black,title={Before you move on}]
What does \ja{ふくらむ} mean? (\textquotedblleft{}To swell; to puff up\textquotedblright{}.) Say it once more.
\end{tcolorbox}
\section[\texorpdfstring{\ja{さます} (samasu)}{さます (samasu)}]{\ja{さます} (samasu) — to let cool}

\label{lesson:JA-C473-samasu}



\begin{quote}
Before the new one: say the Japanese for to scorch, then the Japanese for to swell.
\end{quote}

\subsection*{You'll want to know: \ja{さます}}
\textbf{\ja{さます}} — \emph{samasu} — \textquotedblleft{}to let cool\textquotedblright{}.

\begin{quote}
\textbf{\ja{おちゃを} \ja{すこし} \ja{さまして} \ja{ください。}} — \emph{ocha o sukoshi samashite kudasai} — \textquotedblleft{}Please let the tea cool a little.\textquotedblright{}
\end{quote}

Polite \textbf{\ja{さまします}}, \emph{samashimasu}; a godan verb. Past (polite): \textbf{\ja{ははが} \ja{みそしるを} \ja{さましました。}}, \emph{haha ga misoshiru o samashimashita}.

\begin{cousinweb}
\textbf{\ja{さます}} (\emph{samasu}) is a native Japanese word (wago), inherited rather than borrowed from Chinese; its kanji spelling was attached to it later (JMdict).
\end{cousinweb}

\begin{grammarlens}[title={What you've built}]
One more word for describing what you lived and felt.
\end{grammarlens}

\subsection*{Your turn}
\begin{itemize}
\raggedright
  \item \emph{Say it:} \emph{samasu}
  \item \emph{Say it:} \emph{samasu}, once more
  \item \emph{Say it:} \emph{ocha o sukoshi samashite kudasai} — Please let the tea cool a little.
  \item \emph{Recall:} say the Japanese for to scorch, then the Japanese for to swell, then say \emph{samasu} again
\end{itemize}

\begin{tcolorbox}[breakable,colback=teal!4,colframe=teal!35!black,title={Before you move on}]
What does \ja{さます} mean? (\textquotedblleft{}To let cool\textquotedblright{}.) Say it once more.
\end{tcolorbox}
\section[(dialogue)]{Looking back: from Works of Art to Simmering and Melting}

\label{lesson:JA-R473-review-468-473}



\begin{quote}
Say the Japanese for a work of art and an exhibition.
\end{quote}

\subsection*{Your turn}
Say these aloud:

\begin{itemize}
\raggedright
  \item a work of art, an exhibition, a sculpture, a masterpiece, to assemble
  \item a musical instrument, a small bell, to beat a drum, to play music, the tone of an instrument
  \item a piece of music, sheet music, a folk song, an enka ballad, whistling
  \item a stage, acting, to applaud, to sing in a choir, to compose music
  \item ingredients, a seasoning, a green onion, wheat flour, soybeans
  \item to stew, to melt something, to scorch, to swell, to let cool
\end{itemize}

\begin{tcolorbox}[breakable,colback=teal!4,colframe=teal!35!black,title={Before you move on}]
A work of art? (\textbf{\ja{さくひん}}, \emph{sakuhin}.) An exhibition? (\textbf{\ja{てんらんかい}}, \emph{tenrankai}.) A sculpture? (\textbf{\ja{ちょうこく}}, \emph{chōkoku}.) A masterpiece? (\textbf{\ja{けっさく}}, \emph{kessaku}.) To assemble? (\textbf{\ja{くみたてる}}, \emph{kumitateru}.) A musical instrument? (\textbf{\ja{がっき}}, \emph{gakki}.) A small bell? (\textbf{\ja{すず}}, \emph{suzu}.) To beat a drum? (\textbf{\ja{たたく}}, \emph{tataku}.) To play music? (\textbf{\ja{かなでる}}, \emph{kanaderu}.) The tone of an instrument? (\textbf{\ja{ねいろ}}, \emph{neiro}.) A piece of music? (\textbf{\ja{きょく}}, \emph{kyoku}.) Sheet music? (\textbf{\ja{がくふ}}, \emph{gakufu}.) A folk song? (\textbf{\ja{みんよう}}, \emph{min'yō}.) An enka ballad? (\textbf{\ja{えんか}}, \emph{enka}.) Whistling? (\textbf{\ja{くちぶえ}}, \emph{kuchibue}.) A stage? (\textbf{\ja{ぶたい}}, \emph{butai}.) Acting? (\textbf{\ja{えんぎ}}, \emph{engi}.) To applaud? (\textbf{\ja{はくしゅする}}, \emph{hakushu suru}.) To sing in a choir? (\textbf{\ja{がっしょうする}}, \emph{gasshō suru}.) To compose music? (\textbf{\ja{さっきょくする}}, \emph{sakkyoku suru}.) Ingredients? (\textbf{\ja{ざいりょう}}, \emph{zairyō}.) A seasoning? (\textbf{\ja{ちょうみりょう}}, \emph{chōmiryō}.) A green onion? (\textbf{\ja{ねぎ}}, \emph{negi}.) Wheat flour? (\textbf{\ja{こむぎこ}}, \emph{komugiko}.) Soybeans? (\textbf{\ja{だいず}}, \emph{daizu}.) To stew? (\textbf{\ja{にこむ}}, \emph{nikomu}.) To melt something? (\textbf{\ja{とかす}}, \emph{tokasu}.) To scorch? (\textbf{\ja{こがす}}, \emph{kogasu}.) To swell? (\textbf{\ja{ふくらむ}}, \emph{fukuramu}.) To let cool? (\textbf{\ja{さます}}, \emph{samasu}.)
\end{tcolorbox}
